\chapter{\texttt{org-export-in-background t}: async org export}
\label{ch:org-export}
\minitoc\label{org-export-in-background-t-async-org-export}
\index{org-export-in-background}\index{org-mode!async export}\index{export dispatch}

% ── Opening signpost ──────────────────────────────────────────────
Chapter~\ref{ch:org-modules} trimmed the loading cost of opening an
org file; this chapter tackles the other bottleneck: export.
Exporting a large org file blocks Emacs for the full duration.  This chapter explains \texttt{org-export-in-background}, which
offloads the work to a child \texttt{--batch} process, and covers the
export stack and interaction with per-file export options.

\section{The problem with synchronous export}\label{the-problem-with-synchronous-export}

Exporting a large org file to LaTeX or PDF calls \texttt{org-export-to-file} synchronously, blocking the Emacs UI for the duration. A 200-page document can take 30--120 seconds during which Emacs is completely unresponsive.

\section{How async export works}\label{how-async-export-works}

With \texttt{(setq org-export-in-background t)}, \texttt{org-export-to-file} spawns a \textbf{child Emacs process} in \texttt{-\/-batch} mode. The child:

\begin{enumerate}
\item Loads only what \texttt{org-export-async-init-file} specifies (default: \texttt{$\sim$/.emacs.d/init.el}).
\item Evaluates the export form.
\item Writes the output file to disk and exits.
\end{enumerate}

The parent polls for completion and opens the output file when done. Progress is visible in the \texttt{*Org Export Dispatch*} buffer.

\begin{minted}{elisp}
;; Set in the org use-package :custom block:
(org-export-in-background t)

;; To use a minimal init file for the child process (faster):
(setq org-export-async-init-file
      (expand-file-name "export-init.el" user-emacs-directory))
\end{minted}

\section{\texttt{org-export-stack}}\label{org-export-stack}

During and after async export, \texttt{M-x org-export-stack} shows a buffer listing all export processes with their status (running / done / error) and output files. Keybindings in this buffer allow re-running or opening results.

\section{Interaction with export options}\label{interaction-with-export-options}

The options \texttt{:with-toc}, \texttt{:section-numbers}, \texttt{:with-author}, and \texttt{\#+OPTIONS:} keywords in the org file are honoured by the child process because it loads and evaluates the same org file. Header arguments in the export dispatch (\texttt{C-c C-e}) are passed as the export form's arguments.

% ── Mental model summary ──────────────────────────────────────────
\section{Mental model}
Setting \texttt{org-export-in-background} to \texttt{t} causes the export
dispatcher to spawn a child Emacs in \texttt{--batch} mode, so the parent
Emacs remains responsive during long exports.  The child process loads
\texttt{org-export-async-init-file} instead of the full user init; pointing
this at a minimal init file keeps the child fast and predictable.
\texttt{M-x org-export-stack} provides a dashboard of running and completed
export processes together with their output.  Per-file \texttt{\#+OPTIONS:}
keywords and arguments supplied through the export dispatch (\texttt{C-c C-e})
are honoured by the child process because it loads and evaluates the same org
file.
